\documentclass[11pt]{article}
\usepackage{amsmath,amssymb,amsthm,mathtools}
\usepackage{enumitem}
\usepackage[margin=1.05in]{geometry}
\usepackage[colorlinks=true,linkcolor=black,urlcolor=blue]{hyperref}
\usepackage{booktabs,longtable,array}

\newcommand{\E}{\mathbb E}
\newcommand{\Pp}{\mathbb P}
\newcommand{\Var}{\operatorname{Var}}
\newcommand{\Cov}{\operatorname{Cov}}
\newcommand{\osc}{\operatorname{osc}}
\newcommand{\1}{\mathbf 1}
\newcommand{\cF}{\mathcal F}
\newcommand{\cH}{\mathcal H}
\newcommand{\cG}{\mathcal G}
\newcommand{\cB}{\mathcal B}
\newcommand{\eps}{\varepsilon}
\newcommand{\PD}{\operatorname{PD}(1)}
\newcommand{\Pois}{\operatorname{Poisson}}
\newcommand{\cyc}{\mathrm{cyc}}

\newtheorem{theorem}{Theorem}[section]
\newtheorem{lemma}[theorem]{Lemma}
\newtheorem{proposition}[theorem]{Proposition}
\newtheorem{corollary}[theorem]{Corollary}
\newtheorem{assumption}[theorem]{Assumption}
\theoremstyle{remark}
\newtheorem{remark}[theorem]{Remark}

\title{Audit and rigorous repair of\\
``Mesoscopic cycles of Luce permutations''}
\author{}
\date{18 July 2026}

\begin{document}
\maketitle

\begin{abstract}
The submitted draft does not prove its advertised theorems.  Several errors are
local and repairable, but four are structural: the asserted slab good event is
incompatible with the exponential tail of the clocks; the balance lemma is
used beyond its stopping range and simultaneously over histories without
justification; the history-dependent slab process is treated as a product of
ordinary transition matrices; and the central limit theorem is inferred from
factorial moments with an error far too weak to survive centering.  This note
records the defects precisely and gives a rigorous replacement for the core
orbit argument.  The replacement uses a stopped exploration, adaptive
Chernoff bounds, and oscillation contraction from a two-sided Doeblin bound; it
eliminates the trace argument entirely.  It proves the local cycle law under an
explicit quantitative slab hypothesis.  The expectation and mesoscopic
Poisson conclusions then follow rigorously.  A corrected CLT is also given,
but it requires a quantitative joint local law and a uniform second-moment
bound for short cycles.  What remains unproved is the derivation of the slab
hypothesis from the comparable-weights Luce model, especially the treatment of
late clock slabs.
\end{abstract}

\section{Status of the submitted manuscript}

The principal claims in the abstract and Theorems 1.1--1.5 should not be
presented as proved.  The expectation calculation can be repaired once a
quantitative one-point local law is available.  The mesoscopic Poisson process
can be repaired once a quantitative joint local law is available.  The CLT
needs substantially more than the factorial-moment statement currently
proved.  The claimed reductions for the remaining conjectures are not
well-formed as mathematical statements.

The table below lists the main defects.  Line numbers refer to the submitted
TeX file.

\renewcommand{\arraystretch}{1.15}
\begin{longtable}{p{0.08\textwidth}p{0.25\textwidth}p{0.58\textwidth}}
\toprule
Lines & Issue & Consequence and repair \\
\midrule
\endfirsthead
\toprule
Lines & Issue & Consequence and repair \\
\midrule
\endhead
349--367 & Conditioning contradiction & The sigma-field is said to contain the
sets $H_{n,a}$ and the ordered clocks inside each slab.  Those data determine
the complete ranking and hence the permutation, leaving no random
``within-row matching.''  The environment must contain only the information
actually conditioned on in the cited permanental decomposition. \\
407--421 & Impossible lower slab mass & With $T=(B+1)\log n$, a width-$h$ slab
near $T$ has mass of order $h e^{-\theta T}$, not order $h$.  Therefore
$m_{n,a}\ge cnh$ for every slab, with $c$ independent of $n$, is false.  The
same error invalidates the claimed lower bounds on all cells.  A new tail
construction or a precise external slab lemma is indispensable. \\
529--538 & Wrong Hilbert--Schmidt norm & The equality
$\sum_{a,b}p_bP_{ab}=1$ is false.  On $L^2(p)$ the correct squared
Hilbert--Schmidt norm is
$\sum_{a,b}p_aP_{ab}^2/p_b$, which is bounded by the ratio estimate. \\
556--637 & Overstated balance lemma & The Freedman calculation leaves a term
$(\log m)/\sqrt{m_b}$, and the stopping probability has exponent of order
$\sqrt{m_b}\log m_b$, not
$(\log m_b)^2\sqrt{m_b}/\delta$.  Both corrections are harmless when cell
sizes are polynomial, but the displayed lemma is not proved as stated. \\
640--678 & Circular use of balance & A row-wise lemma valid for at most
$\gamma m_a$ selections is asserted for every global prefix $r\le\delta n$.
A single row may have been selected more than $\gamma m_a$ times.  One must
stop when a row reaches that level and then prove that the stopping time is
unlikely. \\
640--690 & Quantifiers over histories & Freedman's inequality gives a
high-probability assertion along a realized adaptive exploration.  It does not
produce one event on which the estimate holds for every possible history.
The proof must use root-specific stopped events and integrate their failure
probabilities. \\
696--773 & Non-Markov covariance argument & The slab coordinate is not Markov:
the next kernel depends on the residual matching and the whole history.
Different conditional histories do not evolve under a common product of
kernels.  The claimed geometric covariance estimate therefore does not
follow.  It is unnecessary: adaptive Chernoff bounds control row visits using
only the uniform upper ratio $K_t(a,c)\le A p_c$. \\
781--858 & Invalid trace representation & The future random kernels depend on
the path and on the distinguished root.  A product of one realized sequence
of matrices sums over counterfactual paths that would have generated different
kernels.  Thus the displayed trace is not the root average.  The subsequent
run expansion also uses unjustified operator identities.  A uniform
oscillation contraction gives the return density directly and removes this
entire section. \\
839--858 & Incorrect crossover arithmetic & With the stated
$c_*=1/(8|\log q'|)$, the factor
$h^{-1/2}(q')^{k/2}$ at $k=c_*\log n$ is of constant order (up to
$\sqrt{\log n}$), not a negative power of $n$.  Moreover $q'=q_0+\eta_n$ and
$h_k=(q')^k$ define each other circularly. \\
891--924 & False hazard upper bound & For $S=m-1$, the bound obtained in the
proof is $m^{-1}\exp(O(\beta m))$, not $C/m$.  Indeed hazards
$\lambda_j=e^{-\beta}/(m-j)$ give a product of order $m^{-e^{-\beta}}$.
Telescoping is valid only while $S\le\gamma m< m$; the complement must be
controlled probabilistically. \\
1007--1048 & Residual stability not proved & Balanced cell depletion conditional
on the number of visits to a row does not imply that every row is depleted in
the same proportion.  The asserted formulas for $r_a$ and $p^R/p$ require an
additional row-occupancy estimate. \\
1068--1087 & Wrong sampling denominator & Since $V_j$ is an independent uniform
vertex, $\Pp(V_j\in R\mid\mathcal E)=(n-K)/n$, not
$(n-K)/(n-j+1)$.  With the correct denominator the cancellation with the
residual local law is exactly $((n-K)/n)(1/(n-K))=1/n$. \\
460--469 & False macroscopic distribution & The number of $\PD$ parts at least
$\delta$ is not Poisson.  For $\delta>1/2$ it is at most one, whereas a
non-degenerate Poisson variable can exceed one.  Only its expectation
$\log(1/\delta)$ is needed later. \\
1255--1299 & Invalid CLT moment argument & The sum includes fixed lengths
$r>K$, although the local law requires lengths tending to infinity; it also
claims $\ell\delta n=o(n)$ for fixed $\delta$, which is false.  More
fundamentally, relative $o(1)$ errors in raw factorial moments do not survive
the cancellations involved in centered moments. \\
1302--1339 & Tightness does not give variance asymptotics & An $O_p(1)$ remainder
need not have bounded second moment.  Therefore neither the variance claim nor
the covariance error follows from the stated tightness. \\
1347--1388 & Ill-formed mixed estimate & Expressions such as
``$L_{n,j}/n\to s$'' are not finite-$n$ events.  The formula contains an
$o(1)$ inside a limit and confuses vertex-rooted lengths with ranked cycle
lengths.  The process $\sum_v\delta_{L_n(v)/n}$ counts every cycle with
multiplicity equal to its length. \\
1415--1453 & Invalid singular-profile reduction & The comparable-weights local
law cannot be invoked for weights whose ratio grows like $n^\alpha$; the
failure occurs at all scales, not only at $k=\Theta(n)$. \\
1458--1467 & Incorrect scale diagnosis & $n/\log n=o(n)$, and
$c^r=o(1/r)$ there.  Thus the stated reason that the one-point local law is
ineffective is exactly backwards.  The genuine missing input is a sufficiently
uniform joint/extreme-value estimate. \\
\bottomrule
\end{longtable}

\section{A corrected abstract slab framework}

This section isolates a set of hypotheses under which the intended local law
is rigorously true.  The hypotheses are deliberately stated at the level
actually used by the proof.  Establishing them for comparable Luce weights is
a separate theorem; the submitted draft does not do so.

Fix a root $v$.  Let $b_t$ be the slab of the $t$th point in its orbit, and let
$c(v)$ be the row containing the target $v$.  Let rows have sizes
$m_a=np_a$, and let
\[
 P_{ab}=\frac{N_{ab}}{m_a},\qquad \sum_bP_{ab}=1,
 \qquad \sum_a p_aP_{ab}=p_b.
\]

\begin{assumption}[Quantitative slab input]\label{ass:slab}
There are constants $\alpha,\gamma>0$ and $\zeta\in(0,1/2)$, independent of $n$ and of the
scale $k$, such that on an environment event $\cG$:
\begin{enumerate}[label=(H\arabic*)]
\item $p_{\min}:=\min_a p_a\ge n^{-1/2+\zeta}$ and the number of rows is at
most $n^{O(1)}$.
\item The baseline kernel has the two-sided ratio bound
\begin{equation}\label{eq:ratio-correct}
 \alpha p_b\le P_{ab}\le \alpha^{-1}p_b
 \qquad(a,b).
\end{equation}
In particular every cell has size
$N_{ab}\ge \alpha n p_ap_b$.
\item Until a row has lost $\gamma m_a$ targets, every residual target in
that row has conditional selection probability in
$[e^{-\beta}/s,e^\beta/s]$, where $s$ is the residual row size.
The same assertion holds under the root-avoiding law.
\item Under the root-avoiding law, at the $j$th visit to row $c(v)$ the
closure hazard satisfies
\begin{equation}\label{eq:hazard-correct-ass}
 \lambda_j=\frac{e^{\xi_j}}{m_{c(v)}-j},\qquad |\xi_j|\le\beta.
\end{equation}
\item The environment failure probability is $O(n^{-B})$ for every fixed
$B$ after the constants in the construction are chosen.
\end{enumerate}
Finally, the mesh is chosen so that the balance error $\chi_n$ defined in
\eqref{eq:chi} below obeys, for the cycle scale $k$ under consideration,
\begin{equation}\label{eq:parameter-input}
 \beta+\chi_n\le C\bigl(\rho^k+n^{-c}\bigr)
\end{equation}
for some fixed $\rho\in(0,1)$ and $c>0$.
\end{assumption}

The lower bound in (H1) is only a convenient sufficient condition.  What is
really needed below is $np_{\min}\gg\log n$ and polynomially large cells.
The submitted slab construction does not provide either condition near its
clock horizon.

\subsection{Corrected balanced depletion}

\begin{lemma}[Balanced depletion with the correct error]\label{lem:balance-correct}
Let $m$ targets be partitioned into categories of sizes $m_1,\dots,m_d$,
with $m_b\ge m_*\ge m^\vartheta$ for a fixed $\vartheta>0$.  Targets are
selected without replacement in a predictable adaptive manner, and whenever
$s$ targets remain, each individual remaining target has conditional
selection probability in $[e^{-\beta}/s,e^\beta/s]$.  Fix
$\gamma\in(0,1/2)$ and $B<\infty$.  Then, with probability at least
\[
 1-Cd\Bigl(m^{-B}+e^{-c\sqrt{m_*}\log m_*}\Bigr),
\]
simultaneously for all categories $b$ and all $r\le\gamma m$,
\begin{equation}\label{eq:balance-correct}
 \left|X_b(r)-\frac{rm_b}{m}\right|
 \le C_Bm_b\left(
 \beta+\frac{\sqrt{\log m}}{m_b^{1/4}}
       +\frac{\log m}{m_b^{1/2}}
 \right).
\end{equation}
The same estimate, with an additional $O(1/m_b)$ relative error, holds if one
fixed target is forbidden throughout.
\end{lemma}

\begin{proof}
Fix $b$, put $N_r=m_b-X_b(r)$ and $s_r=m-r$, and stop at
$\tau=\inf\{r:N_r<\sqrt{m_b}\}$.  For
$Z_r=\log(N_r/s_r)$, Taylor expansion of the two possible increments gives,
for $r<\tau$,
\[
 \left|\E(Z_{r+1}-Z_r\mid\cF_r)\right|
 \le \frac{C\beta}{s_r}+C\left(\frac1{s_r^2}+\frac1{s_rN_r}\right),
 \qquad
 \E\bigl[(\Delta Z_r-\E_r\Delta Z_r)^2\mid\cF_r\bigr]
 \le\frac{C}{s_rN_r}.
\]
The centered increment is bounded by $C/\sqrt{m_b}$.  Up to time
$\gamma m\wedge\tau$, the predictable quadratic variation is at most
$C/\sqrt{m_b}$ and the total drift is at most
$C(\beta+m_b^{-1/2})$.  Freedman's inequality, with
\[
 x=C_B\left(\frac{\sqrt{\log m}}{m_b^{1/4}}
             +\frac{\log m}{m_b^{1/2}}\right),
\]
gives a failure probability at most $2m^{-B}$.  To hit $\tau$ before
$\gamma m$, the martingale must make a downward excursion of order
$\log m_b$; Freedman's denominator is then of order
$(\log m_b)/\sqrt{m_b}$, yielding the bound
$e^{-c\sqrt{m_b}\log m_b}$.  On the complementary event,
\[
 \left|\log\frac{N_r/m_b}{s_r/m}\right|
 \le C\left(\beta+\frac{\sqrt{\log m}}{m_b^{1/4}}
                   +\frac{\log m}{m_b^{1/2}}\right).
\]
Since $s_r/m\ge1-\gamma$, exponentiation and
$X_b(r)=m_b-N_r$ give \eqref{eq:balance-correct}.  A union bound proves the
simultaneous assertion.  Forbidding one target changes the relevant initial
sizes by at most one and gives the stated additional term.
\end{proof}

Applied row by row, define
\begin{equation}\label{eq:chi}
 \chi_n=C\max_{a,b}\left(
  \beta+\frac{\sqrt{\log n}}{N_{ab}^{1/4}}
       +\frac{\log n}{N_{ab}^{1/2}}
 \right).
\end{equation}
Let $\sigma$ be the first exploration time at which, in some row that has
not yet lost $\gamma m_a$ targets, the estimate
\eqref{eq:balance-correct} fails for one of its cells.  Lemma~\ref{lem:balance-correct},
applied to the row-selection subsequence in each row and then union-bounded,
gives, conditionally on a good environment,
\begin{equation}\label{eq:balance-stop-tail}
 \Pp^\circ(\sigma\le \delta n\mid\cH,\cG)
 \le \varepsilon_{\rm bal,n},
 \qquad
 \varepsilon_{\rm bal,n}
 \le C d^2\left(n^{-B}+e^{-c\sqrt{N_{\min}}\log N_{\min}}\right),
\end{equation}
where $N_{\min}=\min_{a,b}N_{ab}$.  Under (H1)--(H2), the right side is
smaller than any prescribed power of $n$ after increasing $B$.  Introducing
$\sigma$ is essential: one must not condition on the future event that all
balance inequalities hold, since such conditioning would alter the
transition probabilities.

\subsection{Stopping the exploration and controlling row visits}

Let
\[
 \tau=\inf\left\{t:\text{some row $a$ has been used more than }
                         \gamma m_a\text{ times by time }t\right\},
 \qquad T=\tau\wedge\sigma.
\]
Before $T$, summing the one-target bounds over a cell gives a conditional slab
transition kernel $K_t$ satisfying
\begin{equation}\label{eq:kernel-perturb}
 K_t(a,b)=P_{ab}(1+\delta_{t,ab}),\qquad
 |\delta_{t,ab}|\le\eta,
 \qquad \eta=C(\beta+\chi_n).
\end{equation}
This assertion is pathwise only up to $T$; no claim is made that one sample
event controls every counterfactual history.  For estimates of one-time
marginals, continue the slab process with the fixed kernel $P$ after $T$.

\begin{lemma}[Adaptive visit bound]\label{lem:adaptive-visits}
Assume \eqref{eq:ratio-correct} and \eqref{eq:kernel-perturb}.  There is
$\delta_0=\delta_0(\alpha,\gamma)>0$ such that, for $k\le\delta_0 n$,
\begin{equation}\label{eq:visit-tail}
 \Pp^\circ(T\le k\mid\cH,\cG)
 \le \varepsilon_{\rm bal,n}+d\exp(-c n p_{\min}).
\end{equation}
For every fixed row $c$,
\begin{equation}\label{eq:root-row-tail}
 \Pp^\circ\left(S_c(k)>\frac{m_c}{2},\ T>k\mid\cH,\cG\right)
 \le \exp(-c n p_c).
\end{equation}
Consequently the same high-visit probability without $T>k$ is bounded by
the right side of \eqref{eq:root-row-tail} plus the right side of
\eqref{eq:visit-tail}.
\end{lemma}

\begin{proof}
Consider the process continued with $P$ after $T$.  At every step its
conditional probability of visiting row $c$ is at most
\[
 (1+\eta)\alpha^{-1}p_c=:A p_c;
\]
before $T$ this is \eqref{eq:kernel-perturb}, and after $T$ it follows from
\eqref{eq:ratio-correct}.  For adapted Bernoulli variables with conditional
success probabilities at most $Ap_c$, iteration of conditional moment
generating functions gives
\[
 \E\exp\left(\lambda\sum_{t=1}^k\1_{\{b_t=c\}}\right)
 \le \exp\bigl(Akp_c(e^\lambda-1)\bigr).
\]
Choose $\delta_0$ so that $A\delta_0<\gamma/2$.  Chernoff's bound at level
$\gamma np_c$ is then $e^{-cnp_c}$.  Union over rows bounds the event that
$\tau\le k$ occurs before $\sigma$; adding \eqref{eq:balance-stop-tail}
proves \eqref{eq:visit-tail}.  The same calculation with level $np_c/2$
proves \eqref{eq:root-row-tail}, after decreasing $\delta_0$ if necessary.
The possible visit at time zero changes the threshold by one and is absorbed
because $np_{\min}\to\infty$.
\end{proof}

This lemma closes the circularity in the submitted proof: balance is used
only before its own failure time and before row over-depletion, and the
combined stopping time is then shown to be overwhelmingly unlikely.

\subsection{Oscillation mixing instead of the trace argument}

For a probability law $\mu$ on the slab space, write its density relative to
$p$ as $u(a)=\mu(a)/p_a$, and set
$\osc(u)=\max_a u(a)-\min_a u(a)$.

\begin{lemma}[Robust oscillation contraction]\label{lem:oscillation}
Let $P$ be stationary for $p$ and satisfy
$P_{ab}\ge\alpha p_b$.  Let $(\bar K_t)$ be any deterministic sequence of
Markov kernels satisfying
$|\bar K_t(a,b)/P_{ab}-1|\le\eta$, with $\eta\le\alpha/8$.
If $u_t$ is the density relative to $p$ of a chain driven by these kernels,
then, from any point start,
\begin{equation}\label{eq:osc-bound}
 \|u_t-1\|_\infty
 \le C_\alpha(1-\alpha/2)^{t-1}+C_\alpha\eta,
 \qquad t\ge1.
\end{equation}
The same conclusion holds for a history-dependent chain when
$\bar K_t(a,\cdot)$ is defined by averaging the actual next-step law over
histories conditional on the current slab being $a$.
\end{lemma}

\begin{proof}
The density adjoint of $P$ is the reverse kernel
\[
 P^*(b,a)=\frac{p_aP_{ab}}{p_b}\ge\alpha p_a.
\]
Hence $P^*=\alpha\Pi+(1-\alpha)R$, where
$\Pi f=(\E_p f)\1$ and $R$ is Markov.  Therefore
\[
 \osc(P^*f)\le(1-\alpha)\osc(f).
\]
Write $\bar K_t^*=P^*+D_t^*$.  If $u\ge0$ and $\E_pu=1$, then
\[
 |D_t^*u(b)|\le\eta P^*u(b)\le\eta\max u,
\]
so
\[
 \osc(\bar K_t^*u)
 \le(1-\alpha)\osc(u)+2\eta\max u
 \le(1-\alpha+2\eta)\osc(u)+2\eta.
\]
Since $\eta\le\alpha/8$, iteration gives
$\osc(u_t)\le C(1-\alpha/2)^t+C\eta$.  Because $u_t$ is nonnegative with
$p$-mean one, $\|u_t-1\|_\infty\le\osc(u_t)$.  After one step from a point,
\eqref{eq:ratio-correct} and the perturbation bound give a uniformly bounded
density, establishing the initial constant.

For a history-dependent chain, define
\[
 \bar K_t(a,b)=\Pp(b_{t+1}=b\mid b_t=a,\cH),
\]
with an arbitrary admissible row when the conditioning event has probability
zero.  It is a convex combination of admissible conditional kernels and hence
obeys the same entrywise estimate.  The one-time marginals satisfy the usual
recursion with $\bar K_t$, so the deterministic argument applies.  To handle the stopping times, continue the chain with $P$ after $T$; the
continued kernels remain admissible, and the original and continued chains
disagree by at most $\Pp(T\le t)$.
\end{proof}

The key point is that this proof concerns only one-time marginals.  It does not
assert a false Markov property for the full residual process, and it does not
form a path-dependent matrix product or trace.

\subsection{Correct hazard telescoping}

\begin{lemma}[Hazard telescoping away from full depletion]\label{lem:hazard-telescope-correct}
Let $m\ge1$, let $0\le S\le\gamma m$ with fixed $\gamma<1$, and suppose
$\lambda_j=e^{\xi_j}/(m-j)$ with $|\xi_j|\le\beta$.  Then
\begin{equation}\label{eq:hazard-telescope-correct}
 \prod_{j=0}^{S-1}(1-\lambda_j)\lambda_S
 =\frac{1+O_\gamma(\beta)}{m}.
\end{equation}
No bound of the form $C/m$ is asserted uniformly for $S$ close to $m$.
\end{lemma}

\begin{proof}
For the unperturbed hazards $1/(m-j)$, the product is exactly $1/m$.
Divide the perturbed product by the unperturbed one.  The terminal ratio is
$(m-S)\lambda_S=e^{\xi_S}$.  For $j<S$, put $\ell=m-j\ge(1-\gamma)m$.
The mean value theorem gives
\[
 \left|\log(1-e^{\xi_j}/\ell)-\log(1-1/\ell)\right|
 \le C_\gamma\frac{\beta}{\ell}.
\]
Summing and using
$\sum_{j<S}(m-j)^{-1}\le\log(1/(1-\gamma))+O(m^{-1})$ shows that the log of
the total ratio is $O_\gamma(\beta)$.  Exponentiation proves the claim.
\end{proof}

\subsection{The repaired local law}

\begin{theorem}[Local cycle law under the slab input]\label{thm:local-conditional}
Assume Assumption~\ref{ass:slab}.  Then there are constants
$\delta_0,c_0>0$ such that, uniformly for roots $v$ and
$1\le k\le\delta_0n$,
\begin{equation}\label{eq:local-conditional}
 \Pp(L_n(v)=k+1)
 =\frac1n\left[1+O\left((1-\alpha/2)^k+\beta+\chi_n\right)\right]
 +O(n^{-1-c_0}).
\end{equation}
Consequently, if \eqref{eq:parameter-input} holds, then for every
$r_n\to\infty$,
\[
 \Pp(L_n(I_n)=r)=\frac1n\left(1+O(\widetilde\rho^{\,r})+O(n^{-c_0})\right)
\]
uniformly for $r_n\le r\le\delta_0n$, for some
$\widetilde\rho\in(0,1)$.
\end{theorem}

\begin{proof}
Work first conditional on a good environment.  Continue the root-avoiding
slab process with $P$ after the stopping time $T$.  By
Lemma~\ref{lem:oscillation}, its time-$k$ density satisfies
\[
 u_k(c(v))=1+O\bigl((1-\alpha/2)^k+\eta\bigr).
\]
By Lemma~\ref{lem:adaptive-visits}, replacing the continued process by the
original process changes the density at $c(v)$ by at most
$[\varepsilon_{\rm bal,n}+d e^{-cnp_{\min}}]/p_{c(v)}$, which is smaller than
every prescribed power of $n$ under (H1)--(H2).  The exact root-avoiding hazard identity is
\[
 \Pp(L_n(v)=k+1\mid\cH)
 =\E^\circ\left[
  \1_{\{b_k=c(v)\}}
  \prod_{j=0}^{S(k)-1}(1-\lambda_j)\lambda_{S(k)}
 \right].
\]
On $S(k)\le m_{c(v)}/2$, Lemma~\ref{lem:hazard-telescope-correct} replaces
the hazard product by $(1+O(\beta))/m_{c(v)}$.  On the complement, the integrand is at most one.  Lemma~\ref{lem:adaptive-visits}gives probability at most
$e^{-cnp_{c(v)}}+\varepsilon_{\rm bal,n}+d e^{-cnp_{\min}}$.  Therefore
\[
 \Pp(L_n(v)=k+1\mid\cH)
 =\frac{1+O(\beta)}{m_{c(v)}}
   \Pp^\circ(b_k=c(v)\mid\cH)
  +O\!\left(e^{-cnp_{c(v)}}+\varepsilon_{\rm bal,n}
                 +d e^{-cnp_{\min}}\right).
\]
Since $m_{c(v)}=np_{c(v)}$, the main term is
$n^{-1}(1+O(\beta))u_k(c(v))$.  Substituting the oscillation estimate and
$\eta=C(\beta+\chi_n)$ proves the conditional bound.  The balance-stopping and environment failure probabilities may be made
$O(n^{-B})$ with $B>1+c_0$, and are then absorbed after removing the
conditioning.  Finally use
\eqref{eq:parameter-input} and enlarge the geometric rate.
\end{proof}

\begin{remark}[What remains to prove for Luce permutations]
Theorem~\ref{thm:local-conditional} repairs the probabilistic argument, but it
does not validate the submitted Proposition 3.1.  To obtain the advertised
unconditional Luce theorem, one still needs a correct construction of clock
slabs satisfying Assumption~\ref{ass:slab}, including the late exponential
tail.  The assertion that all fixed-width slabs up to $(B+1)\log n$ have mass
of order $h$ cannot be used.
\end{remark}

\section{Correct downstream consequences}

\subsection{Expectation}

The expectation argument is valid after two corrections: use the quantitative
bound of Theorem~\ref{thm:local-conditional}, and replace the false Poisson
claim for the number of macroscopic $\PD$ parts by its expectation.

\begin{proposition}[Expectation from three inputs]\label{prop:expectation-correct}
Assume:
\begin{enumerate}[label=(E\arabic*)]
\item for every fixed $r$, $\E C_r(\sigma_n)\to\lambda_r$, and
$\sum_{r\ge1}|\lambda_r-r^{-1}|<\infty$;
\item for some $q\in(0,1)$ and $c>0$,
\[
 \left|\E C_r(\sigma_n)-\frac1r\right|
 \le\frac{C}{r}(q^r+n^{-c})
 \quad (r\le\delta_0n);
\]
\item for every fixed $\delta>0$,
\[
 \E\#\{i:L_{n,i}>\delta n\}\longrightarrow\log(1/\delta).
\]
\end{enumerate}
Then
\[
 \E N_n^{\cyc}=\log n+\gamma_{\rm E}
 +\sum_{r\ge1}\left(\lambda_r-\frac1r\right)+o(1).
\]
\end{proposition}

\begin{proof}
Fix $K$ and $\delta\le\delta_0$.  Split at $K$ and $\delta n$.  The fixed
part tends to $\sum_{r\le K}\lambda_r$.  The quantitative local law gives
\[
 \sum_{K<r\le\delta n}\E C_r
 =\sum_{K<r\le\delta n}\frac1r+O(q^K)+o(1).
\]
The macroscopic input contributes $\log(1/\delta)+o(1)$.  Since
$\sum_{K<r\le\delta n}r^{-1}=\log n+\log\delta-\log K+o(1)$, the
$\log\delta$ terms cancel.  Let $n\to\infty$, then $K\to\infty$, and use
$H_K-\log K\to\gamma_{\rm E}$ and absolute summability.
\end{proof}

\subsection{Mesoscopic Poisson process}

For fixed-dimensional mesoscopic point-process convergence, relative
$o(1)$ errors are sufficient because the limiting factorial moments are
finite.  The correct joint input is as follows.

\begin{assumption}[Quantitative joint rooted local law]\label{ass:joint}
For each fixed $\ell$ there are $q\in(0,1)$ and $c>0$ such that, uniformly for
lengths $r_1,\dots,r_\ell$ with $\min_jr_j\to\infty$ and
$\sum_jr_j=o(n)$,
\begin{equation}\label{eq:joint-correct}
 \Pp\left(
  L_n(V_j)=r_j\ \forall j,
  \text{ and the rooted cycles are distinct}
 \right)
 =\frac1{n^\ell}\left[
 1+O_\ell\left(\sum_{j=1}^\ell q^{r_j}+n^{-c}\right)
 \right].
\end{equation}
\end{assumption}

The residual exploration proof can yield this statement if the quantitative
slab input is hereditary after deletion and the row-stopping estimates are
made conditional with failure probabilities smaller than $n^{-B}$ for
arbitrarily large $B$.  In the induction step, the correct identity is
\[
 \frac{n-K}{n}\cdot\frac1{n-K}=\frac1n.
\]

\begin{proposition}[Mesoscopic Poisson process]\label{prop:ppp-correct}
Under Assumption~\ref{ass:joint}, for $a_n\to\infty$, $a_n=o(n)$, the point
process of cycle lengths rescaled by $a_n$ converges vaguely on
$(0,\infty)$ to the Poisson point process with intensity $dx/x$.
\end{proposition}

\begin{proof}
For disjoint compact intervals $J_j=[u_j,v_j]\subset(0,\infty)$ and fixed
nonnegative integers $\ell_j$, expand the joint falling factorial of the cycle
counts.  Marking each selected cycle by a vertex gives the exact identity
\[
 \E\mathfrak C_n(r_1,\dots,r_\ell)
 =\frac{n^\ell}{\prod_jr_j}
 \Pp(\text{the event in \eqref{eq:joint-correct}}).
\]
The error in \eqref{eq:joint-correct} is uniform over all tuples in the
windows.  Hence the mixed factorial moments converge to
\[
 \prod_j\left(\sum_{r\in a_nJ_j}\frac1r\right)^{\ell_j}
 \longrightarrow
 \prod_j\left(\log\frac{v_j}{u_j}\right)^{\ell_j}.
\]
These are the mixed factorial moments of independent Poisson variables.
Tightness on a compact subset follows from the uniformly bounded first moment,
and the standard factorial-moment criterion gives vague convergence.
\end{proof}

\subsection{A CLT that the local laws can actually support}

The submitted fixed-$K$, fixed-$\delta$ proof cannot be repaired by changing a
line or two.  A valid route uses moving cutoffs and an error that is
polynomially small in $n$, not merely relative $o(1)$.

Set
\[
 a_n=\lceil(\log n)^2\rceil,
 \qquad b_n=\left\lfloor\frac{n}{\sqrt{\log n}}\right\rfloor,
 \qquad
 M_n=\sum_{a_n<r\le b_n}C_r,
 \qquad
 \mu_n=\sum_{a_n<r\le b_n}\frac1r.
\]
Then $\mu_n\sim\log n$, while any fixed number of selected lengths in this
range has total $o(n)$.

\begin{assumption}[Short-cycle second moment]\label{ass:short2}
Uniformly for $x\le(\log n)^2$,
\begin{equation}\label{eq:short2}
 \E\left(\sum_{r\le x}C_r\right)^2\le C(1+\log x)^2.
\end{equation}
\end{assumption}

A sufficient condition for \eqref{eq:short2} is the uniform pair bound
$\E C_rC_s\le C/(rs)$ off the diagonal and
$\E(C_r)_2\le C/r^2$, together with $\E C_r\le C/r$ in the relevant range
and finite-length domination.  Fixed-dimensional convergence alone does not
supply this uniform estimate.

\begin{theorem}[Corrected CLT criterion]\label{thm:clt-correct}
Assume Assumptions~\ref{ass:joint} and \ref{ass:short2}, with the error in
\eqref{eq:joint-correct} valid quantitatively for all
$a_n<r_j\le b_n$.  Assume also the one-point local bound through
$\delta_0n$ and a uniformly bounded expected number of cycles longer than
$\delta_0n$.  Then
\[
 \frac{N_n^{\cyc}-\E N_n^{\cyc}}{\sqrt{\log n}}
 \Longrightarrow N(0,1),
 \qquad
 \Var N_n^{\cyc}=\log n+o(\log n).
\]
\end{theorem}

\begin{proof}
For fixed $\ell$, expand the falling factorial of $M_n$.  By
\eqref{eq:joint-correct},
\[
 \E(M_n)_\ell
 =\mu_n^\ell+O_\ell\left(
   \mu_n^{\ell-1}\sum_{r>a_n}\frac{q^r}{r}
   +n^{-c}\mu_n^\ell
 \right)
 =\mu_n^\ell+O_A((\log n)^{-A})
\]
for every fixed $A$, after increasing $n$.  Thus, for each fixed moment order,
the factorial moments agree with those of $Z_n\sim\Pois(\mu_n)$ to an error
smaller than every negative power of $\mu_n$.  Passing from factorial to
ordinary moments and then centering is now legitimate: the polynomial factors
introduced by centering are powers of $\mu_n$, which are dominated by the
error above.  Therefore
\[
 \E\left(\frac{M_n-\mu_n}{\sqrt{\mu_n}}\right)^j
 -\E\left(\frac{Z_n-\mu_n}{\sqrt{\mu_n}}\right)^j\longrightarrow0
\]
for each fixed $j$.  Standardized Poisson moments converge to Gaussian
moments, and the Gaussian is moment-determinate.  Hence
$(M_n-\mu_n)/\sqrt{\mu_n}\Rightarrow N(0,1)$ and
$\Var M_n=\mu_n+o(\mu_n)$.

Let
\[
 R_n^{\rm short}=\sum_{r\le a_n}C_r,
 \qquad
 R_n^{\rm long}=\sum_{r>b_n}C_r.
\]
Assumption~\ref{ass:short2} gives
$\E(R_n^{\rm short})^2=O((\log\log n)^2)=o(\log n)$.  For the long part,
the one-point local law on $(b_n,\delta_0n]$ and the macroscopic expectation
bound give
\[
 \E R_n^{\rm long}=O(\log\log n).
\]
Deterministically $R_n^{\rm long}\le n/b_n\le2\sqrt{\log n}$, so
\[
 \E(R_n^{\rm long})^2
 \le 2\sqrt{\log n}\,\E R_n^{\rm long}
 =o(\log n).
\]
Consequently $R_n=R_n^{\rm short}+R_n^{\rm long}$ satisfies
$\E R_n^2=o(\log n)$, and therefore
$(R_n-\E R_n)/\sqrt{\log n}\to0$ in $L^2$.  Since
$\mu_n\sim\log n$, Slutsky's theorem gives the CLT for the total number of
cycles.  Finally,
\[
 \Var(N_n^{\cyc})=\Var(M_n)+o(\log n)
\]
by Cauchy--Schwarz and the $L^2$ bound on $R_n$, proving the variance
asymptotic.
\end{proof}

\section{Correct formulations for the remaining conjectures}

\subsection{Mixed mesoscopic--macroscopic independence}

A rigorous mixed estimate should be stated in terms of factorial moments and
bounded continuous test functions.  Let
$\Xi_n=\sum_i\delta_{L_{n,i}/a_n}$ count cycles once, and let
$\mathbf P_n=(L_{n,1}/n,L_{n,2}/n,\dots)$.  A sufficient statement is:
for every bounded continuous $G$ on $S^\downarrow$, every family of disjoint
compact intervals $J_1,\dots,J_m\subset(0,\infty)$, and fixed
$\ell_1,\dots,\ell_m$,
\begin{equation}\label{eq:mixed-correct}
 \E\left[G(\mathbf P_n)
       \prod_{j=1}^m(\Xi_n(J_j))_{\ell_j}\right]
 \longrightarrow
 \E G(\mathbf P)\,
 \prod_{j=1}^m\left(\int_{J_j}\frac{dx}{x}\right)^{\ell_j},
\end{equation}
where $\mathbf P\sim\PD$.  Together with tightness, this identifies the joint
limit and its independence.  Equation \eqref{eq:mixed-correct} is a genuine
missing theorem; it is not implied by separate marginal convergence.

\subsection{Endpoint-singular profiles}

No consequence for weights with ratio $\kappa_n\to\infty$ follows from a
theorem whose constants depend on a fixed comparable-weights parameter
$\kappa$.  A valid reduction would have to assume, or prove, an analogue of
Assumption~\ref{ass:slab} with constants and errors controlled for the
singular profile.  Without that input, the proposed Proposition 9.3 should be
removed.

\subsection{The scale \texorpdfstring{$n/\log n$}{n/log n}}

The scale $n/\log n$ is submacroscopic.  The one-point error
$q^r+n^{-c}$ is excellent there.  An extreme-gap theorem instead needs joint
control over a moving family of rare long-cycle events, or a direct coupling
of the upper-tail point process.  The corrected obstruction should be stated
in those terms, not as a failure of the one-point local law.

\section{Recommended revision of the original paper}

A mathematically honest revision should make the following changes.

\begin{enumerate}[label=\arabic*.]
\item Change the title and abstract from a claimed resolution to a conditional
reduction unless a correct tail-slab theorem is supplied.
\item Replace the current good-event proposition by an exact statement of the
slab result actually available from the cited manuscript.  Do not assert
uniform $m_a\asymp nh$ for fixed-width slabs up to order $\log n$.
\item Replace Sections 4.3--4.4 (mixing, covariance, and trace averaging) by
Lemmas~\ref{lem:adaptive-visits} and \ref{lem:oscillation} above.
\item Replace the global hazard upper bound by
Lemma~\ref{lem:hazard-telescope-correct} and bound the high-visit event
directly by adaptive Chernoff.
\item State residual stability only after proving a stopped row-occupancy
bound, and correct the independent-root denominator from $n-j+1$ to $n$.
\item Correct the macroscopic proposition: retain convergence to $\PD$ and the
expectation $\log(1/\delta)$, but delete the Poisson assertion.
\item Keep the expectation theorem, conditional on the repaired local law.
Keep the mesoscopic Poisson theorem, conditional on the quantitative joint
law.
\item Replace the CLT section by Theorem~\ref{thm:clt-correct}, and explicitly
prove or assume the uniform short-cycle second-moment bound.
\item Replace the three ``remaining conjecture'' subsections by the corrected
formulations in the preceding section.
\item Relabel the numerical appendix as numerical illustration.  Simulation
cannot verify a theorem, and the displayed variance ratios around $0.85$ do
not by themselves establish convergence to one.
\end{enumerate}

\section{Bottom line}

The expectation and mesoscopic Poisson calculations are close to correct once
a genuine quantitative local and joint local law is supplied.  The submitted
proof of that local law is not correct, but its main probabilistic mechanism
can be repaired by stopping and oscillation contraction, as shown above.  The
remaining unresolved mathematical task is concrete: construct a slab
decomposition for comparable exponential rates that satisfies
Assumption~\ref{ass:slab} uniformly across the late clock tail, and prove its
heredity under finitely many residual cycle explorations.  Until that is done,
the advertised resolutions of the conjectures should remain conditional.

\end{document}